% ============================================================
%  APPENDIX A -- KEY CODE IMPLEMENTATIONS
%  AI MockPrep Project Report
% ============================================================

\chapter{Key Code Implementations}
\label{app:code}

This appendix presents representative source-code excerpts that underpin the
core functionality of AI MockPrep.  All listings are in TypeScript (the
project-wide language) unless noted otherwise.

% ────────────────────────────────────────────────────────────
\section{Session Authentication Helper}
\label{app:auth}
% ────────────────────────────────────────────────────────────

The \texttt{getCurrentUser} helper (\texttt{lib/auth.ts}) is called from every
server component and server action that needs the identity of the signed-in
user.  It reads the \texttt{session} cookie set by Firebase Auth, verifies it
with the Admin SDK, and merges Auth-profile data with the richer profile stored
in Firestore.

\begin{lstlisting}[language=Java,
  caption={Session verification and user-profile resolution (\texttt{lib/auth.ts})},
  label={lst:getCurrentUser}]
import { cookies }  from "next/headers";
import { redirect } from "next/navigation";
import { auth, db } from "@/firebase/admin";

export async function getCurrentUser() {
  const cookieStore  = await cookies();
  const sessionCookie = cookieStore.get("session");
  if (!sessionCookie) return null;

  try {
    const decoded    = await auth.verifySessionCookie(
                         sessionCookie.value, true);
    const userRecord = await auth.getUser(decoded.uid);

    let photoURL    = userRecord.photoURL;
    let displayName = userRecord.displayName;
    let bio = "", headline = "";

    if (db) {
      const doc = await db.collection("users")
                          .doc(decoded.uid).get();
      if (doc.exists) {
        const d = doc.data()!;
        photoURL    = d.avatar     || d.photoURL    || photoURL;
        displayName = d.name       || d.displayName || displayName;
        bio         = d.bio        || "";
        headline    = d.headline   || "";
      }
    }

    return {
      uid: decoded.uid, email: decoded.email,
      name: displayName || decoded.name || "User",
      photoURL: photoURL || undefined,
      bio, headline,
    };
  } catch {
    return null;
  }
}

export async function requireAuth() {
  const user = await getCurrentUser();
  if (!user) redirect("/sign-in");
  return user;
}
\end{lstlisting}

% ────────────────────────────────────────────────────────────
\section{Interview Creation Server Action}
\label{app:interview-create}
% ────────────────────────────────────────────────────────────

The \texttt{createInterview} server action (\texttt{lib/actions/interview.action.ts})
validates the caller's identity, delegates question generation to the
\texttt{questionGenerator} module, and persists the new interview document to
Firestore.

\begin{lstlisting}[language=Java,
  caption={Interview creation server action
           (\texttt{lib/actions/interview.action.ts})},
  label={lst:createInterview}]
"use server";
import { db }               from "@/firebase/admin";
import { generateQuestions } from "@/lib/interview/questionGenerator";
import { getCurrentUser }    from "@/lib/auth";

export async function createInterview(params: {
  userId: string; role: string; level: string;
  techstack: string[]; type: string; company?: string;
}) {
  const { role, level, techstack, type, company } = params;
  const currentUser = await getCurrentUser();
  if (!currentUser)
    return { success: false, message: "Not authenticated" };

  const questions = await generateQuestions(
    role, level, techstack, type);

  const data = {
    userId: currentUser.uid,
    company: company || null,
    role, level, techstack, type, questions,
    finalized: false,
    createdAt: new Date().toISOString(),
  };

  const docRef = await db.collection("interviews").add(data);
  return { success: true, interviewId: docRef.id };
}
\end{lstlisting}

% ────────────────────────────────────────────────────────────
\section{Question Generator}
\label{app:qgen}
% ────────────────────────────────────────────────────────────

Question generation first consults a curated CSV bank; if no rows match the
role/level/tech-stack query it falls back to a generative model
(\texttt{lib/interview/questionGenerator.ts}).

\begin{lstlisting}[language=Java,
  caption={Question-generation entry point
           (\texttt{lib/interview/questionGenerator.ts})},
  label={lst:questionGenerator}]
"use server";
import {
  loadQuestionsFromCSV,
  generateFallbackCSVQuestions,
} from "./csvQuestions";

export const generateQuestions = async (
  role: string, level: string,
  techstack: string[], type = "Technical"
): Promise<string[]> => {

  const csvQuestions = await loadQuestionsFromCSV(
    role, level, techstack, type);

  if (csvQuestions.length > 0) return csvQuestions;

  console.warn("No CSV match -- using fallback generation");
  return generateFallbackCSVQuestions(role, level, techstack);
};
\end{lstlisting}

% ────────────────────────────────────────────────────────────
\section{Speech Coach -- Per-Response Analysis}
\label{app:speech-coach}
% ────────────────────────────────────────────────────────────

\texttt{analyzeSpeechCoachResponse} (\texttt{lib/speech/coach.ts}) takes a
single voice-transcribed answer and produces a multi-dimensional score covering
confidence, tone, clarity, pacing, and filler-word control.

\begin{lstlisting}[language=Java,
  caption={Per-response speech-coach analysis
           (\texttt{lib/speech/coach.ts})},
  label={lst:speechCoach}]
export const analyzeSpeechCoachResponse = (
  question: string, answer: string,
  durationSeconds: number
): SpeechCoachResponseInsight => {

  const words      = answer.trim().split(/\s+/).filter(Boolean);
  const wordCount  = words.length;
  const wpm        = Math.round((wordCount / durationSeconds) * 60);

  // Filler-word detection
  const fillerHits  = countFillerHits(answer);
  const fillerCount = Object.values(fillerHits)
    .reduce((s, c) => s + c, 0);
  const fillerPer100 = (fillerCount / wordCount) * 100;

  // Pattern-match scoring
  const assertiveCount =
    countPatternMatches(answer, ASSERTIVE_PATTERNS);
  const hedgingCount   =
    countPatternMatches(answer, HEDGING_PATTERNS);
  const structureScore =
    STRUCTURE_CUES.filter(c =>
      answer.toLowerCase().includes(c)).length;

  const confidence   = clamp(60 + assertiveCount*10
                                - hedgingCount*8);
  const pacing       = clamp(100 - Math.abs(wpm-145)*1.1);
  const fillerControl= clamp(100 - fillerPer100*9);
  const clarity      = clamp(40 + structureScore*10);

  const overall = clamp(
    confidence*0.24 + tone*0.16 +
    clarity*0.24   + pacing*0.20 +
    fillerControl*0.16
  );

  return {
    question, answer, durationSeconds,
    confidence, tone, clarity, pacing,
    fillerControl, overall,
    fillerCount, fillerHits, wpm,
  };
};
\end{lstlisting}

% ────────────────────────────────────────────────────────────
\section{Speech Coach -- Session Summary Aggregation}
\label{app:speech-summary}
% ────────────────────────────────────────────────────────────

After all questions are answered, \texttt{summarizeSpeechCoachResponses}
aggregates per-question insights into a session-level summary with actionable
recommendations.

\begin{lstlisting}[language=Java,
  caption={Session-level speech-coach summary
           (\texttt{lib/speech/coach.ts})},
  label={lst:speechSummary}]
export const summarizeSpeechCoachResponses = (
  analyses: SpeechCoachResponseInsight[]
): InterviewSpeechCoachSummary => {
  if (analyses.length === 0)
    return { analyzedResponses: 0, /* ... zero values */ };

  const average = (fn: (a: SpeechCoachResponseInsight)
                        => number) =>
    Math.round(analyses.reduce((s, a) => s + fn(a), 0)
               / analyses.length);

  // Aggregate filler-word frequency map
  const fillerMap: Record<string, number> = {};
  for (const a of analyses)
    for (const [word, n] of Object.entries(a.fillerHits))
      fillerMap[word] = (fillerMap[word] ?? 0) + n;

  const topFillerWords = Object.entries(fillerMap)
    .sort((a, b) => b[1] - a[1])
    .slice(0, 3).map(([w]) => w);

  const summary = {
    analyzedResponses: analyses.length,
    confidence:   average(a => a.confidence),
    tone:         average(a => a.tone),
    clarity:      average(a => a.clarity),
    pacing:       average(a => a.pacing),
    fillerControl:average(a => a.fillerControl),
    overall:      average(a => a.overall),
    totalFillerWords:
      analyses.reduce((s, a) => s + a.fillerCount, 0),
    topFillerWords,
    averageWpm: average(a => a.wpm),
    insights: [], recommendations: [],
  };

  // Build human-readable insights
  if (summary.confidence >= 75)
    summary.insights.push(
      "Strong confidence language in your spoken responses.");
  if (summary.fillerControl < 70)
    summary.recommendations.push(
      "Replace filler words with a brief pause to keep "
      + "delivery crisp.");

  return summary;
};
\end{lstlisting}

% ────────────────────────────────────────────────────────────
\section{Inadequate Response Detection}
\label{app:inadequate}
% ────────────────────────────────────────────────────────────

Before passing an answer to the AI evaluator, the interview engine screens for
non-substantive responses.  Detected inadequate answers receive a fixed low
score rather than consuming API quota on empty content.

\begin{lstlisting}[language=Java,
  caption={Inadequate response detection
           (\texttt{lib/realTimeAnalysis.ts})},
  label={lst:inadequate}]
private isInadequateResponse(answer: string): boolean {
  const clean = answer.toLowerCase().trim();

  const dontKnow = [
    /^i don'?t know$/, /^don'?t know$/,
    /^no idea$/, /^not sure$/,
    /^idk$/, /^pass$/, /^skip$/,
  ];
  if (dontKnow.some(p => p.test(clean))) return true;

  // Reject very short answers (< 10 chars)
  if (clean.length < 10) return true;

  // Reject pure filler / yes-no answers
  const filler = [
    /^(um+|uh+|well|so|like|you know)[\s\.,]*$/,
    /^(yes|no|maybe|possibly|probably)[\s\.,]*$/,
  ];
  if (filler.some(p => p.test(clean))) return true;

  return false;
}
\end{lstlisting}

% ────────────────────────────────────────────────────────────
\section{Route-Protection Middleware}
\label{app:middleware}
% ────────────────────────────────────────────────────────────

The \texttt{proxy} function (\texttt{proxy.ts}) runs on every request via
Next.js Middleware.  It allows public paths unconditionally and redirects
unauthenticated users to \texttt{/sign-in}, preserving the intended destination
in a \texttt{next} query parameter.

\begin{lstlisting}[language=Java,
  caption={Route-protection middleware (\texttt{proxy.ts})},
  label={lst:middleware}]
const PUBLIC_PATHS = new Set([
  "/", "/sign-in", "/sign-up", "/verify-email",
  "/about", "/contact", "/help", "/privacy", "/terms",
]);

export async function proxy(request: NextRequest) {
  const { pathname } = request.nextUrl;

  // Skip Next.js internals and static assets
  if (pathname.startsWith("/api/") ||
      pathname.startsWith("/_next/") ||
      pathname.match(/\.(svg|png|jpg|jpeg|ico|js|css)$/))
    return NextResponse.next();

  if (PUBLIC_PATHS.has(pathname))
    return NextResponse.next();

  const session = request.cookies.get("session")?.value;
  if (!session || !auth)
    return redirectToSignIn(request, pathname);

  try {
    await auth.verifySessionCookie(session, true);
    return NextResponse.next();
  } catch {
    return redirectToSignIn(request, pathname);
  }
}
\end{lstlisting}
